% Seminário 2/2026 — Engenharia Econômica para Software (PUC Minas)
% Tema 1: Estimativa baseada em histórias de usuário
%
% Estrutura: capa -> sumário -> seções (artigo/secoes/) -> referências (referencias.bib).
% Compilar com: python build.py   (na raiz do repositório)

\documentclass[article,12pt,a4paper,oneside,brazil,section=TITLE]{abntex2}

\usepackage[T1]{fontenc}
\usepackage[utf8]{inputenc}
\usepackage{mathptmx}
\usepackage{microtype}
\usepackage{indentfirst}
\usepackage{graphicx}
\usepackage{amsmath}
\usepackage{booktabs}
\usepackage{array}
\usepackage[alf]{abntex2cite}

\hypersetup{
  pdftitle   = {Estimativa baseada em histórias de usuário},
  pdfauthor  = {Carlos José Gomes Batista Figueiredo; Bernardo Souza Alvim; Mateus Araujo Santos; Gabriela Alvarenga Cardoso, Marcos Alberto Ferreira Pinto, Rafael Ganascini de Moura},
  pdfsubject = {Seminário 2/2026 — Engenharia Econômica para Software — PUC Minas},
  colorlinks = true,
  linkcolor  = black,
  citecolor  = black,
  urlcolor   = black
}

% Tipografia ABNT: Times em todo o texto; títulos de seção em negrito, mesmo corpo do texto.
\renewcommand{\ABNTEXchapterfont}{\rmfamily\bfseries}
\renewcommand{\ABNTEXsectionfontsize}{\normalsize}
\renewcommand{\ABNTEXsubsectionfontsize}{\normalsize}

% Quadros (ABNT): flutuante próprio, numerado separadamente das tabelas.
\newfloat{quadro}{loq}{Quadro}

\OnehalfSpacing

\begin{document}

\pagestyle{empty}
% Capa — NBR 14724 (adaptada ao enunciado do seminário).
% Para incluir ou corrigir integrantes, edite apenas a lista abaixo.
% Logotipo opcional: salve em artigo/figuras/logo-pucminas.png e ele entra sozinho.

\newcommand{\integrante}[2]{%
  \begin{minipage}[t]{0.48\textwidth}\centering
    #1\\{\footnotesize #2}
  \end{minipage}}

\begingroup
\thispagestyle{empty}
\SingleSpacing
\centering

\IfFileExists{figuras/logo-pucminas.png}{%
  \includegraphics[height=2.2cm]{figuras/logo-pucminas.png}\par\vspace{0.5cm}}{}
{\bfseries PONTIFÍCIA UNIVERSIDADE CATÓLICA DE MINAS GERAIS\par}
Instituto de Ciências Exatas e Informática\par
Curso de Graduação em Engenharia de Software\par

\vspace{1.4cm}

\integrante{Carlos José Gomes Batista Figueiredo}{1507022@pucminas.edu.br}\hfill
\integrante{Bernardo Souza Alvim}{1528669@pucminas.edu.br}\\[1.6ex]
\integrante{Mateus Araujo Santos}{1487920@pucminas.edu.br}\hfill
\integrante{Gabriela Alvarenga Cardoso}{1026227@pucminas.edu.br}\\[1.6ex]
\integrante{Marcos Alberto Ferreira Pinto}{1523749@pucminas.edu.br}\hfill
\integrante{Rafael Ganascini de Moura}{1456345@pucminas.edu.br}

\vfill

{\Large\bfseries ESTIMATIVA BASEADA EM HISTÓRIAS DE USUÁRIO\par}
\vspace{0.4cm}
{\large Story points, \textit{velocity} e previsões financeiras em projetos ágeis\par}

\vfill

\hfill\begin{minipage}{0.52\textwidth}
  \small\raggedright
  Artigo do Seminário 2/2026 apresentado à disciplina Engenharia Econômica
  para Software, 6º período, campus Lourdes.\par
  \vspace{0.3cm}
  Docente: Prof. Michelle Hanne Soares de Andrade
\end{minipage}

\vfill

Belo Horizonte\par
27 de outubro de 2026\par
\endgroup
\newpage


\renewcommand{\contentsname}{SUMÁRIO}
\renewcommand{\bibname}{REFERÊNCIAS}
\tableofcontents*
\newpage

\pagestyle{simple}
\section{Introdução}
\label{sec:introducao}

\subsection{Contextualização}

A Engenharia Econômica para Software trata das decisões de investimento em
sistemas sob incerteza: quanto custa construir, em quanto tempo e se o retorno
esperado justifica o desembolso. Todas essas perguntas dependem de uma
estimativa de esforço, já que o esforço, valorado pelo custo do tempo da
equipe, é o principal componente do custo de um projeto de software
\cite{boehm1981}. Não surpreende que a estimativa de custos seja um dos temas
mais estudados da área: \citeonline{jorgensen2007} identificaram 304 artigos
sobre o assunto em 76 periódicos.

Nas últimas duas décadas, boa parte desse planejamento migrou de planos
detalhados e preditivos para ciclos curtos e adaptativos. Desde o Manifesto
Ágil \cite{agilemanifesto2001}, métodos como Scrum e Extreme Programming (XP)
se difundiram: o 17º \textit{State of Agile Report}, pesquisa com 788
profissionais, registra que 71\% dos respondentes usam práticas ágeis no ciclo
de desenvolvimento de software \cite{digitalai2024}. Nesses métodos, os
requisitos são escritos como histórias de usuário, técnica originária do XP
\cite{beck2004,cohn2004}, e estimados em \textit{story points}, unidades
relativas de tamanho e complexidade, e não em horas ou em reais
\cite{cohn2005}. A soma dos pontos concluídos em uma iteração define a
\textit{velocity} da equipe, que serve de base para prever quantas iterações
faltam para esvaziar o \textit{backlog} \cite{cohn2005}. O Guia do Scrum, por
sua vez, não prescreve uma técnica de dimensionamento e atribui essa
responsabilidade aos desenvolvedores \cite{scrumguide2020}.

\subsection{Problema e relevância econômica}

Como o \textit{story point} não é uma unidade monetária, é preciso uma ponte
entre a escala de pontos e a escala financeira. A ponte mais simples combina
três grandezas: o \textit{backlog} restante $B$ (em pontos), a \textit{velocity}
média $\bar{v}$ (pontos por \textit{sprint}) e o custo da equipe por
\textit{sprint} $c_s$. Com elas obtêm-se o número previsto de \textit{sprints}
$N$, o custo por ponto $c_p$ e o custo total previsto $C$:
%
\begin{equation}
  N = \frac{B}{\bar{v}}, \qquad
  c_p = \frac{c_s}{\bar{v}}, \qquad
  C = N\,c_s = B\,c_p .
  \label{eq:previsao}
\end{equation}
%
Por exemplo, com valores hipotéticos de $B = 240$ pontos,
$\bar{v} = 30$ pontos por \textit{sprint} e $c_s = \text{R\$}\,60.000$,
tem-se $N = 8$ \textit{sprints}, $c_p = \text{R\$}\,2.000$ por ponto e
$C = \text{R\$}\,480.000$.

A simplicidade da Equação~\ref{eq:previsao} esconde o problema econômico do
tema: cada termo herda a incerteza da estimativa que o produz. O
\textit{story point} é a métrica de tamanho mais usada em projetos ágeis
\cite{usman2014,fernandezdiego2020}, e sua atribuição costuma ser feita por
julgamento de especialistas, em especial com \textit{Planning Poker}
\cite{fernandezdiego2020}. Há indícios, porém, de que os pontos atribuídos por
humanos nem sempre refletem o tempo de desenvolvimento efetivamente gasto: em
37.440 histórias de 37 projetos de código aberto, \citeonline{tawosi2022}
encontram evidências de que o \textit{story point} pode não ser um bom
indicador de esforço. A
\textit{velocity}, por sua vez, não é constante, e por isso há trabalhos
dedicados a prevê-la \cite{hearty2009} e a antecipar se uma iteração em curso
vai entregar o que foi planejado \cite{choetkiertikul2018}. Por fim,
\citeonline{jorgensen2005} mostra que as avaliações de incerteza em custo de
software tendem ao excesso de confiança, ou seja, a intervalos mais estreitos
do que a realidade justifica. Erros nesses termos chegam ao orçamento, ao
cronograma e à decisão de investir.

A pesquisa recente tem tentado reduzir esse erro com aprendizado de máquina.
\citeonline{choetkiertikul2019} propõem o Deep-SE, que estima \textit{story
points} a partir do texto das \textit{issues} e superou as referências de
comparação em 23.313 \textit{issues} de 16 projetos de código aberto;
\citeonline{fu2023} seguem a mesma linha com um modelo baseado em
transformadores (GPT2SP). Entretanto, ao replicar o Deep-SE,
\citeonline{tawosi2023} constatam que ele supera as estimativas de mediana e o
TF-IDF com SVM em poucos casos com significância estatística. Esse debate
reforça a importância de comparar qualquer estimador com um \textit{baseline}
bem definido, prática que, para a estimativa de esforço em geral, ajuda a
reduzir a instabilidade das conclusões \cite{sarro2018}. Do lado gerencial, o uso das métricas ágeis em
escala também é desafiador: \citeonline{schull2023} mostram, em estudo de caso
numa grande montadora, a dificuldade de reportar progresso e desempenho de
equipes ágeis aos níveis de produto e portfólio. E, na ótica da economia da
informação, \citeonline{koch2008} ilustra, para projetos de código aberto, que
modelos de esforço baseados na participação de programadores estimam um esforço
sensivelmente menor do que modelos baseados em medidas de saída, como linhas de
código, isto é, a escolha da medida de esforço altera as conclusões sobre a
eficiência econômica do desenvolvimento.

\subsection{Objetivos}

O objetivo geral deste trabalho é analisar, com base na literatura, como as
estimativas baseadas em histórias de usuário (\textit{story points} e
\textit{velocity}) sustentam previsões de esforço, prazo e custo em projetos
ágeis, e quais são os limites dessas previsões. Constituem objetivos
específicos:
%
\begin{alineas}
  \item conceituar histórias de usuário, \textit{story points} e
        \textit{velocity};
  \item descrever como a produtividade da equipe converte pontos em prazo e em
        custo;
  \item examinar as evidências empíricas sobre a acurácia dessas estimativas;
  \item apresentar exemplos e aplicações, acadêmicos e de mercado; e
  \item discutir limitações e boas práticas para o uso financeiro das
        estimativas.
\end{alineas}

\subsection{Procedimentos metodológicos}

Trata-se de uma pesquisa bibliográfica, exploratória e qualitativa. Partiu-se
dos periódicos sugeridos no enunciado do seminário, a saber, \textit{IEEE
Transactions on Software Engineering} (TSE), \textit{ACM Transactions on
Software Engineering and Methodology} (TOSEM), \textit{Information Systems and
e-Business Management} (ISeBM) e \textit{Information Economics and Policy}
(IEP), complementados por revisões sistemáticas, anais de eventos e livros de
referência da área. Os dados bibliográficos (volume, número, páginas e DOI)
foram conferidos no Crossref. Nas buscas realizadas, os estudos dedicados a
\textit{story points} e \textit{velocity} concentraram-se no TSE; do TOSEM, do
ISeBM e do IEP foram aproveitados trabalhos adjacentes, sobre
\textit{baselines} de estimativa, reporte de métricas em organizações ágeis e
modelagem econômica de esforço, respectivamente. Essa distribuição desigual é
uma limitação do levantamento e não esgota a literatura disponível.

\subsection{Organização do trabalho}

Além desta introdução, o artigo está organizado em três seções. A
Seção~\ref{sec:marco} apresenta o marco teórico sobre histórias de usuário,
\textit{story points}, \textit{velocity} e sua relação com previsões
financeiras. A Seção~\ref{sec:exemplos} traz exemplos e aplicações, e a
Seção~\ref{sec:conclusao} reúne as conclusões. Seguem-se as referências
bibliográficas.

\section{Marco teórico}
\label{sec:marco}

% Marco teórico, parte 1 (pessoa 4): histórias de usuário, story points, Planning Poker
% e comparação com outras técnicas de estimativa.
% A parte 2 (velocity, produtividade e conversão em custo) fica em outro arquivo
% (02b-...), incluído em 02-marco-teorico.tex, para evitar conflitos de merge.

\subsection{Histórias de usuário}
\label{sub:historias}

Uma história de usuário descreve, em poucas linhas, uma funcionalidade do ponto
de vista de quem a usa ou dela se beneficia. A técnica é originária do XP
\cite{beck2004} e foi detalhada por \citeonline{cohn2004}. Ron Jeffries a
resumiu em três elementos: o cartão (\textit{card}), com a descrição curta da
história; a conversa (\textit{conversation}), em que o entendimento da equipe e
do cliente se aprofunda; e a confirmação (\textit{confirmation}), que define
como saber que a história está concluída \cite{mountaingoatstories}. Um modelo
comum de redação é ``Como [tipo de usuário], preciso [fazer algo], para [obter
um benefício]'' (tradução livre de modelo apresentado em
\citeonline{mountaingoatstories}).

Para a estimativa, a qualidade da história importa tanto quanto o método.
\citeonline{wake2003} propõe o acrônimo INVEST: uma boa história é independente
(\textit{independent}), negociável (\textit{negotiable}), valiosa
(\textit{valuable}), estimável (\textit{estimable}), pequena (\textit{small}) e
testável (\textit{testable}). Dois desses critérios incidem diretamente sobre o
tema deste artigo. Ser estimável significa ser compreendida o bastante para
receber uma estimativa aproximada, que ajude o cliente a priorizá-la e
agendá-la. Ser pequena significa, segundo o autor, exigir em geral no máximo
algumas semanas-pessoa de trabalho. Quando a história é grande demais ou falta
informação para estimá-la, o passo seguinte pode ser detalhá-la, fazer um
estudo exploratório (\textit{spike}), dividi-la ou conversar com as partes
interessadas \cite{mountaingoatpoker}.

\subsection{\textit{Story points}}
\label{sub:storypoints}

O \textit{story point} é a unidade com que as histórias são dimensionadas. Ele
expressa o esforço necessário para concluí-las em escala relativa: uma história
de dois pontos deve exigir cerca do dobro do esforço de uma de um ponto, e o que
importa são as razões entre os valores, não os números em si
\cite{mountaingoatpoints}. A mesma fonte indica que a estimativa deve refletir
a quantidade de trabalho, a complexidade e o risco ou a incerteza da história.
Como a escala é construída pela própria equipe, o mesmo número pode representar
esforços diferentes em equipes diferentes. O dimensionamento costuma ser feito
por analogia: a equipe escolhe histórias de referência e compara as demais com
elas \cite{cohn2005}.

Os valores atribuídos vêm, em geral, de uma escala discreta inspirada na
sequência de Fibonacci. O baralho de \textit{Planning Poker} da Mountain Goat
Software, por exemplo, usa 0, 1, 2, 3, 5, 8, 13, 20, 40 e 100
\cite{mountaingoatpoker}, e arredonda alguns valores, como 20 no lugar de 21,
para não sugerir uma precisão que a estimativa não tem \cite{mountaingoatfib}.

Como o ponto não é uma unidade de tempo nem de dinheiro, ele só se converte em
prazo e custo por meio da \textit{velocity} da equipe, tratada na sequência
deste marco teórico. Além disso, a premissa de que os pontos refletem o esforço
real tem apoio empírico limitado, como visto na Seção~\ref{sec:introducao}
\cite{tawosi2022}.

\subsection{\textit{Planning Poker}}
\label{sub:planningpoker}

O \textit{Planning Poker} é uma técnica de estimativa em grupo, baseada no
julgamento de especialistas, usada para atribuir pontos às histórias. A
estimativa por julgamento de especialistas é objeto de ampla revisão em
\citeonline{jorgensen2004}, e o \textit{Planning Poker} foi proposto por
\citeonline{grenning2002} como uma forma de estruturá-la em grupo. No
procedimento original, cada programador escolhe uma carta em segredo, as cartas
são reveladas ao mesmo tempo e, se os valores divergem, quem deu a estimativa
mais alta e quem deu a mais baixa explicam seus motivos, e o grupo discute até
chegar a um consenso. As cartas traziam os valores 1, 2, 3, 5, 7, 10 e infinito,
em dias ideais de programação ou em pontos sem unidade, e o infinito indicava
uma história grande demais, que devia ser dividida \cite{grenning2002}.

Na versão hoje difundida, o \textit{Product Owner} apresenta o item e responde
às dúvidas, mas não estima; os desenvolvedores escolhem as cartas em privado e
as revelam juntos; havendo divergência, os estimadores de valor mais alto e mais
baixo expõem primeiro suas razões, e a rodada se repete até que haja
entendimento suficiente para registrar um valor \cite{mountaingoatpoker}. Cartas
especiais indicam que o item não foi compreendido o bastante para ser estimado
ou que ainda é grande demais. A mesma fonte desaconselha tirar a média das
estimativas, pois ela esconde as razões da divergência.

A revelação simultânea busca evitar a ancoragem, em que uma estimativa dita
primeiro puxa as demais em sua direção \cite{mountaingoatpoker}; o efeito da
ancoragem na estimativa de software é estudado por \citeonline{aranda2005}. A
literatura empírica compara o \textit{Planning Poker} a outras formas de chegar
a estimativas. \citeonline{haugen2006} o contrapõe à estimativa em grupo não
estruturada em uma equipe de XP; \citeonline{molokken2008}, em estudos
industriais, investigam se combinar estimativas por consenso reduz o otimismo
das estimativas individuais; e \citeonline{mahnic2012} tratam do uso da técnica
na estimativa de histórias de usuário. A revelação simultânea e a discussão são
mecanismos para reduzir vieses individuais, não para eliminar a incerteza da
estimativa, que continua se propagando para a previsão de prazo e custo
(Equação~\ref{eq:previsao}).

\subsection{Comparação com outras técnicas de estimativa}
\label{sub:comparacao}

Estimar esforço e custo em projetos ágeis não se resume a \textit{story
points}. O Quadro~\ref{quad:tecnicas} compara-os a duas outras abordagens: a
estimativa direta em horas ou dias ideais e a contagem de pontos de função. A
primeira é a mais próxima do raciocínio financeiro, já que o tempo é o que se
paga; o próprio \citeonline{grenning2002} admite cartas em dias ideais ou em
pontos sem unidade. A segunda mede o tamanho do software pelas funções que ele
oferece ao usuário, e não pelo esforço de construí-las, e foi relacionada ao
esforço de desenvolvimento e às linhas de código por \citeonline{albrecht1983}.

\begin{quadro}[htbp]
  \caption{Comparação entre três formas de estimar o tamanho ou o esforço de
  uma funcionalidade}
  \label{quad:tecnicas}
  \small
  \begin{tabular}{@{}p{2.2cm}p{4.0cm}p{3.9cm}p{4.4cm}@{}}
    \toprule
    \textbf{Técnica} & \textbf{O que se estima} & \textbf{Vantagem} & \textbf{Limitação} \\
    \midrule
    \raggedright Horas ou dias ideais
      & \raggedright Tempo de trabalho, estimado por julgamento de especialistas.
      & \raggedright Converte-se diretamente em prazo e custo.
      & \raggedright O valor varia com quem executa a tarefa e tende a sugerir mais precisão do que a estimativa tem.
      \tabularnewline
    \addlinespace
    \raggedright \textit{Story points}
      & \raggedright Esforço relativo entre histórias, em escala própria da equipe, definido por consenso (por exemplo, no \textit{Planning Poker}).
      & \raggedright Dispensa converter a história em tempo no momento da estimativa; a conversão vem depois, pela \textit{velocity}.
      & \raggedright Escala não comparável entre equipes; exige iterações concluídas para virar prazo e custo, e sua relação com o esforço real é discutível \cite{tawosi2022}.
      \tabularnewline
    \addlinespace
    \raggedright Pontos de função
      & \raggedright Tamanho funcional, obtido pela contagem das funções oferecidas ao usuário (entradas, saídas, consultas e arquivos).
      & \raggedright Mede o produto, não quem o constrói, e não depende da linguagem de programação.
      & \raggedright Exige requisitos descritos com algum detalhe e treino para a contagem; a produtividade (esforço por ponto) ainda precisa ser calibrada.
      \tabularnewline
    \bottomrule
  \end{tabular}
  \par\vspace{1ex}
  \fonte{Elaborado pelos autores, com base em \citeonline{cohn2005}, \citeonline{grenning2002}, \citeonline{albrecht1983} e \citeonline{mountaingoatpoints}.}
\end{quadro}

Em comum, as três abordagens produzem uma medida de tamanho ou de esforço que
só se torna prazo e custo depois de dividida pela produtividade: o termo $B$ da
Equação~\ref{eq:previsao} muda de unidade (horas, pontos ou pontos de função),
mas a lógica da previsão permanece. O que diferencia as técnicas é quanto
conhecimento elas exigem no momento da estimativa e quanto da conversão para
prazo e custo fica para depois.


% Pessoa 5 (velocity, produtividade e conversão de estimativa em custo):
% criar secoes/02b-velocity-custo.tex e incluir aqui com % Marco teórico, parte 2 (pessoa 5): velocity, produtividade e conversão de estimativa em custo.
% A parte 1 (histórias de usuário, story points e Planning Poker) fica em 02a-historias-pontos.tex.
% Os valores numéricos desta parte são hipotéticos e seguem o exemplo da Equação (eq:previsao).
% Convenção desta parte: sem travessões e sem ponto e vírgula. Citações com vários autores
% usam \citeonline, porque \cite separa os sobrenomes com ponto e vírgula no padrão ABNT.

\subsection{\textit{Velocity}}
\label{sub:velocity}

Saber o tamanho de uma história não diz, sozinho, quando ela vai ficar pronta.
Quem responde a essa pergunta é a \textit{velocity}, isto é, a quantidade de
trabalho que a equipe consegue concluir em uma iteração. Quando as histórias são
estimadas em \textit{story points}, a \textit{velocity} é simplesmente a soma dos
pontos das histórias terminadas no \textit{sprint} \cite{cohn2005}. A regra de
contagem é rígida de propósito. Uma história que ficou 90\% pronta vale zero
pontos naquele \textit{sprint}, porque ainda vai consumir esforço no seguinte, e
contá-la pela metade daria a impressão de um avanço que não aconteceu
\cite{cohn2005}.

Há um efeito curioso nessa medida. Se a equipe subestima as histórias de forma
sistemática, ela também termina menos pontos por iteração do que esperava, e a
divisão entre as duas grandezas continua dando um prazo coerente.
\citeonline{cohn2005} vê nisso uma forma de autocorreção da estimativa. Mas a
compensação depende de duas condições que nem sempre se mantêm. A primeira é que
a equipe use a escala de pontos sempre do mesmo jeito. A segunda é que a própria
equipe não mude muito. Basta trocar parte dos integrantes, adotar uma tecnologia
nova ou rever o critério de pontuação para que o histórico deixe de servir de
referência.

Na prática, a \textit{velocity} oscila de um \textit{sprint} para outro. Férias,
feriados, incidentes em produção e histórias que não cabem na iteração bastam
para tirá-la do lugar. Por isso, \citeonline{cohn2005} sugere trabalhar com a
média das últimas iterações e comunicar prazos como uma faixa, e não como um
número fechado. Para uma equipe sem histórico, o autor aponta três caminhos.
Ela pode usar dados de um projeto parecido, rodar algumas iterações antes de
prometer qualquer prazo ou estimar a capacidade disponível quebrando as
histórias em tarefas medidas em horas. Se $v_1, \dots, v_k$ são as
\textit{velocities} das $k$ últimas iterações, a média e os extremos observados
são
%
\begin{equation}
  \bar{v} = \frac{1}{k}\sum_{i=1}^{k} v_i, \qquad
  v_{\min} = \min_i v_i, \qquad
  v_{\max} = \max_i v_i .
  \label{eq:velocity}
\end{equation}

Prever a \textit{velocity} virou, inclusive, tema de pesquisa.
\citeonline{hearty2009} representam a \textit{velocity} de projetos XP com uma
rede bayesiana dinâmica, que se recalibra a cada iteração e, assim, melhora a
previsão à medida que o projeto gera dados. \citeonline{choetkiertikul2018}
seguem outro caminho e usam informações de iterações e de \textit{issues} de
projetos de código aberto para prever, com o \textit{sprint} ainda em andamento,
se a equipe vai entregar o que planejou. Os dois trabalhos partem da mesma
constatação. Tratar a \textit{velocity} como um número fixo esconde boa parte do
risco da previsão.

\subsection{Produtividade e capacidade da equipe}
\label{sub:produtividade}

Em economia, produtividade é a razão entre o que se produz e o que se gasta para
produzir. Na estimativa tradicional de software, ela costuma aparecer como
tamanho por esforço, em linhas de código por pessoa-mês, como nos modelos de
\citeonline{boehm1981}, ou em pontos de função por unidade de esforço, medida
que \citeonline{albrecht1983} relacionaram ao esforço de desenvolvimento. A
\textit{velocity} tem a mesma forma, pontos por iteração, mas carrega duas
diferenças que mudam a maneira de usá-la.

A primeira diferença está na unidade. Como vimos na
Seção~\ref{sub:storypoints}, o ponto de uma equipe não equivale ao ponto de
outra. Comparar a \textit{velocity} de duas equipes, portanto, é comparar
réguas diferentes, e \citeonline{cohn2005} desaconselha esse uso. O problema
piora quando a \textit{velocity} vira meta de desempenho. A equipe passa a
``entregar'' mais pontos sem entregar mais funcionalidade, bastando para isso
aumentar as estimativas, e a métrica perde justamente o que tinha de útil, que
era ajudar a prever. O estudo de caso de \citeonline{schull2023}, numa grande
montadora, mostra como é difícil levar indicadores de equipes ágeis aos níveis
de produto e portfólio sem distorcê-los. Já \citeonline{koch2008} mostra, em
projetos de código aberto, que a escolha da medida de esforço altera as
conclusões sobre a eficiência econômica do desenvolvimento.

A segunda diferença é que a \textit{velocity} mistura produtividade com
capacidade. Se ela cai, pode ser que a equipe tenha rendido menos, ou apenas que
tenha trabalhado menos dias. Uma forma simples de separar os dois efeitos é
corrigir a média pela capacidade prevista. Sendo $h_p$ as horas-pessoa
disponíveis no próximo \textit{sprint} e $\bar{h}$ a média de horas-pessoa das
iterações usadas no cálculo de $\bar{v}$, a \textit{velocity} esperada é
%
\begin{equation}
  v_{\text{esp}} = \bar{v}\,\frac{h_p}{\bar{h}} .
  \label{eq:capacidade}
\end{equation}
%
Uma equipe que costuma entregar 30 pontos e terá 20\% menos horas por causa de
um feriado e de uma licença deveria planejar algo perto de 24 pontos para o
\textit{sprint}. O ajuste supõe que a produtividade por hora não mude, o que é
razoável no curto prazo. Ela deixa de valer quando entra gente nova no time, porque as
horas disponíveis aumentam ao mesmo tempo em que o rendimento por hora cai,
enquanto os novos integrantes aprendem o sistema.

\subsection{Conversão da estimativa em prazo e custo}
\label{sub:custo}

A Equação~\ref{eq:previsao} já transforma o \textit{backlog} em número de
\textit{sprints} e em custo. Para que esse resultado sustente uma decisão de
investimento, ainda faltam duas coisas. É preciso saber de onde vem o custo da
equipe por \textit{sprint} e é preciso dizer o quanto a previsão pode errar.

O custo por \textit{sprint} vem do tempo das pessoas, que é o maior componente do
custo de um projeto de software \cite{boehm1981}. Com $n$ integrantes dedicados,
$h$ horas de trabalho de cada um no \textit{sprint} e um custo por hora $c_h$ que
já inclua salário, encargos e benefícios, tem-se
%
\begin{equation}
  c_s = n\,h\,c_h .
  \label{eq:custosprint}
\end{equation}
%
Os R\$\,60.000 por \textit{sprint} usados na Introdução correspondem, por
exemplo, a seis pessoas trabalhando 80 horas cada em um \textit{sprint} de duas
semanas, a R\$\,125 por hora. Gastos que não dependem do tamanho do time, como
licenças e infraestrutura, entram em $c_s$ quando crescem com a duração do
projeto, ou são somados uma única vez ao custo total quando não crescem.

Como a \textit{velocity} oscila, faz mais sentido prever uma faixa do que um
valor. Usando os extremos da Equação~\ref{eq:velocity} e arredondando o número de
\textit{sprints} para cima, já que ninguém paga meio \textit{sprint} de uma
equipe dedicada, chega-se a um cenário otimista e a um pessimista:
%
\begin{equation}
  N_{\min} = \left\lceil \frac{B}{v_{\max}} \right\rceil, \qquad
  N_{\max} = \left\lceil \frac{B}{v_{\min}} \right\rceil, \qquad
  N_{\min}\,c_s \le C \le N_{\max}\,c_s .
  \label{eq:intervalo}
\end{equation}
%
Suponha que as cinco últimas \textit{velocities} da equipe do exemplo tenham sido
26, 31, 28, 34 e 31 pontos. A média é exatamente 30, o mesmo valor usado na
Introdução, mas o \textit{backlog} de 240 pontos pode levar de 8 a 10
\textit{sprints}, e o custo fica entre R\$\,480.000 e R\$\,600.000. Os
R\$\,480.000 calculados antes são o melhor caso dessa faixa. Apresentar só esse
número a quem aprova o orçamento seria cair no excesso de confiança descrito por
\citeonline{jorgensen2005}. Vale lembrar que os extremos de cinco iterações não
são limites garantidos. Modelos que produzem uma distribuição de prazos, como a
rede bayesiana de \citeonline{hearty2009}, permitem dizer também qual é a
probabilidade de cada cenário.

A incerteza também muda ao longo do projeto. No começo, com requisitos ainda
vagos, ela é grande, e vai diminuindo conforme o produto é entendido e
construído. Esse padrão ficou conhecido como cone da incerteza \cite{boehm1981}.
\citeonline{little2006}, analisando dados de projetos reais, confirma que a
incerteza nas estimativas de prazo é alta, mas encontra um cone menos estreito
do que o modelo sugere e uma razão entre prazo real e prazo estimado com
distribuição assimétrica, de cauda longa para o lado dos atrasos. Para quem faz
a análise financeira, a lição é dupla. O cenário pessimista merece mais peso que
o otimista, e a previsão precisa ser refeita a cada \textit{sprint}, com o
\textit{backlog} e a \textit{velocity} atualizados.

\subsection{Acompanhamento do custo durante o projeto}
\label{sub:agileevm}

Depois que o orçamento é aprovado, as mesmas grandezas servem para acompanhar a
execução. \citeonline{sulaiman2006} adaptaram ao Scrum a gestão de valor
agregado (\textit{Earned Value Management}, EVM), técnica clássica de controle
de projetos, e chamaram o resultado de AgileEVM. A ideia central é medir o
progresso pelos pontos concluídos, sem depender de um cronograma detalhado de
tarefas. Considere $BAC$ o orçamento total (\textit{budget at completion}), $N$ o
número de \textit{sprints} planejados, $j$ os \textit{sprints} já concluídos,
$P_c$ os pontos entregues até agora e $AC$ o custo real gasto até o momento
(\textit{actual cost}). Com essas grandezas calculam-se o valor planejado $PV$,
o valor agregado $EV$ e os índices de desempenho de custo ($CPI$) e de prazo
($SPI$):
%
\begin{equation}
  PV = \frac{j}{N}\,BAC, \qquad
  EV = \frac{P_c}{B}\,BAC, \qquad
  CPI = \frac{EV}{AC}, \qquad
  SPI = \frac{EV}{PV} .
  \label{eq:agileevm}
\end{equation}

Voltando ao exemplo, com $BAC = \text{R\$}\,480.000$ e $N = 8$, imagine que, ao
fim do quarto \textit{sprint}, a equipe tenha entregado 112 dos 240 pontos e
gastado R\$\,250.000. Nesse caso, $PV = \text{R\$}\,240.000$,
$EV = \text{R\$}\,224.000$, $CPI \approx 0{,}90$ e $SPI \approx 0{,}93$. Em
outras palavras, o projeto entregou menos do que o planejado para a metade do caminho
e gastou mais do que vale o que entregou. Se esse ritmo continuar, o custo
estimado no término é $BAC/CPI \approx \text{R\$}\,536.000$. O valor cai dentro
da faixa da Equação~\ref{eq:intervalo}, mas fica bem acima dos R\$\,480.000 que
teriam sido prometidos se a previsão fosse pontual. O AgileEVM não resolve a
incerteza dos pontos, já que $EV$ depende de $B$, e $B$ muda sempre que o
\textit{backlog} muda. O que ele oferece é uma tradução. O progresso, que o time
enxerga em pontos, passa a ser descrito em custo e prazo, que é a linguagem de
quem decide sobre o investimento.

Olhando o caminho inteiro, da história de usuário até o orçamento, há três
conversões. A equipe estima o tamanho em pontos, a \textit{velocity}
transforma pontos em tempo e o custo do time transforma tempo em dinheiro. Cada
uma delas soma incerteza à anterior. Por isso, uma previsão financeira baseada
em \textit{story points} só ajuda de verdade quando chega a quem decide como
uma faixa de valores, revista a cada iteração, e não como um número único
tratado como certo.


\section{Exemplos e aplicações}
\label{sec:exemplos}

\textbf{[Seção em elaboração pelo grupo.]}

% Conclusão (pessoa 6). Convenção igual à do marco teórico 2: sem travessões e sem ponto e vírgula.

\section{Conclusão}
\label{sec:conclusao}

Este trabalho analisou, com base na literatura, como as estimativas por histórias
de usuário sustentam previsões de esforço, prazo e custo em projetos ágeis. A
resposta curta é que elas sustentam, com ressalvas. O \textit{story point} não é
uma unidade monetária e só vira prazo e custo por duas conversões em cadeia. A
\textit{velocity} transforma pontos em tempo, e o custo da equipe transforma
tempo em dinheiro (Seção~\ref{sub:custo}). A Equação~\ref{eq:previsao} é simples
e útil, mas cada um de seus termos herda a incerteza de quem o produziu.

Retomando os objetivos específicos, o marco teórico conceituou histórias de
usuário, \textit{story points} e \textit{velocity} e mostrou como a produtividade
da equipe liga pontos a prazo e a custo. As evidências empíricas reunidas na
Introdução e na Seção~\ref{sec:exemplos} são menos favoráveis do que a prática
sugere. A relação entre pontos e esforço real tem apoio limitado
\cite{tawosi2022}, os modelos de aprendizado de máquina nem sempre superam
referências simples \cite{tawosi2023}, e as avaliações de incerteza tendem ao
excesso de confiança \cite{jorgensen2005}. Os exemplos hipotéticos mostraram
duas aplicações econômicas. Uma faixa de 8 a 10 \textit{sprints} manteve o VPL
positivo, de modo que a incerteza não mudou a decisão de investir, e o custo por
ponto permitiu ordenar funcionalidades pelo retorno, como faz o WSJF.

Daí saem algumas boas práticas para o uso financeiro das estimativas:
%
\begin{alineas}
  \item apresentar o custo como uma faixa, calculada com os extremos da
        \textit{velocity}, e não como um número único (Equação~\ref{eq:intervalo});
  \item dar mais peso ao cenário pessimista, já que a distribuição dos atrasos
        tem cauda longa \cite{little2006};
  \item refazer a previsão a cada \textit{sprint} e acompanhar o custo e o prazo
        com indicadores como $CPI$ e $SPI$ \cite{sulaiman2006};
  \item não comparar pontos nem \textit{velocity} entre equipes e não usar a
        \textit{velocity} como meta de desempenho, o que a faria perder o valor
        preditivo; e
  \item comparar qualquer estimador, humano ou automático, com um
        \textit{baseline} simples, como a mediana histórica \cite{sarro2018}.
\end{alineas}

O trabalho tem limitações. É uma pesquisa bibliográfica e exploratória, com
amostra de fontes desigual entre os periódicos consultados, como já dito nos
procedimentos metodológicos. Os exemplos são hipotéticos e usam valores
escolhidos para ilustrar o cálculo, sem dados de um projeto real. As fontes de
mercado são documentação de fornecedores e não medem a acurácia das previsões.
Esses limites indicam caminhos para estudos futuros, como aplicar o método a
dados históricos de uma organização, para verificar quão larga precisa ser a
faixa de custo, e comparar a previsão por pontos com a previsão em horas ou em
pontos de função em projetos de mesmo tipo.

Em síntese, estimar por histórias de usuário continua sendo uma forma prática e
barata de planejar. O cuidado necessário é não tratar o número que sai da
conta como certo. O que o gestor deve levar à decisão de investimento é uma
faixa de valores, com seus pressupostos explícitos, revista a cada iteração.


\bibliography{referencias}

% Apêndices e anexos (se houver): \appendix e \input{secoes/apendice-a}

\end{document}
